\section{The Dataflow Abstraction}
\label{slab:sec:dataflow}
\newcounter{mycounter}

The introduction argued that a rule-free search for rewrites needs a
signal: something computable from the input query alone that
distinguishes candidates worth examining from the vastly more numerous
candidates that are not. This section develops that signal. It first
makes precise the intuition that equivalent queries are related through
what they compute (\cref{slab:sec:dataflow:insight}), then defines
dataflows as the formal carrier of that intuition
(\cref{slab:sec:dataflow:language}), and finally gives a procedure that
extracts the dataflows of any query in our language
(\cref{slab:sec:dataflow:extraction}). The next section describes how
the synthesis algorithm puts the extracted dataflows to work.

\subsection{Dataflows as a Semantic Signal}
\label{slab:sec:dataflow:insight}

Consider what survives when a query is rewritten. The syntax may change
beyond recognition: subqueries appear or dissolve, a self-join becomes
a window function, a \sqlhavingcolor clause becomes a \sqlwherecolor
over a derived table. What persists is the computation. If the original
query sums a column and filters on the sum, then any equivalent query
must, in some guise, still obtain that sum; the result depends on it.
The essential computations of a query, the ways in which data from the
input tables is combined, transformed, and compared on the way to the
output, are properties of the query's meaning, not of its spelling. We
call each such essential computation a \emph{dataflow}: a path that
information takes from input tables through the query's operations.

This observation cuts in both directions, and both matter for search.
In one direction, a candidate query that performs the input query's
computations is a plausible rewrite: since the input query is correct,
its dataflows are, by construction, sufficient raw material for
producing the correct result. In the other direction, a candidate that
performs computations foreign to the input query is suspect. If the
input query sums a column, a candidate that only averages it is
unlikely to produce the same answers; if the two queries filter on
unrelated columns, equivalence is improbable. Foreign dataflows are
evidence against equivalence, and a search can use that evidence to
set such candidates aside.

One asymmetry deserves emphasis: a rewrite need not use \emph{all} of
the input query's dataflows. Input queries, especially poorly written
ones, often perform computations their results do not depend on; a
redundant join is the classic case. A good rewrite discards exactly
this redundant work, and its dataflows are then a strict subset of the
input's. The signal we want is therefore one-sided. Sharing the
input's dataflows is encouraged; exhibiting dataflows the input lacks
is penalized; failing to use everything the input computes is not
penalized at all, because that failure is often precisely the
optimization.

Finally, note what kind of signal this is. Dataflows are semantic:
they are indifferent to how a computation is written, and in
particular to how large or how syntactically elaborate a candidate is.
A search guided by dataflows can therefore rank a structurally
ambitious candidate above a superficially simple one whenever the
former better preserves the input's computations. This is what allows
the approach to find deep restructurings early, rather than only after
exhausting all syntactically smaller candidates, and it is the
property that distinguishes dataflow guidance from the size-ordered
enumeration that is standard in program synthesis.

\subsection{The Dataflow Language}
\label{slab:sec:dataflow:language}

We now define dataflows formally.

\begin{definition}[Dataflow]
A dataflow $\dataflowchain$ (defined in Figure~\ref{slab:fig:dataflowlanguage}) for a given query $\query$ is a sequence of SQL operations that are performed during $\query$'s execution on one or more input tables or their columns.
\end{definition}

The grammar in \cref{slab:fig:dataflowlanguage} builds dataflows from
the ground up. In the base cases, a dataflow is an input table
$\dbtable$ or a column of an input table. Base dataflows are
deliberately anchored to the \emph{input} tables rather than to
intermediate results: a column alias manufactured halfway through a
query is an artifact of that query's particular phrasing, and tying
dataflows to it would make them syntax-dependent, defeating their
purpose. The recursive cases then record operations applied on top of
existing dataflows. An aggregate application
$\emph{agg}(\dataflowchain_1, \mydots, \dataflowchain_n)$ records that
an aggregate function is computed over the data described by its
argument dataflows; window functions are recorded analogously. The
composition $[\dataflowchain_1, \mydots, \dataflowchain_n]$ records
the assembly of columns into a table, as a \sqlselectcolor clause
does; because composition captures it, \sqlselectcolor itself needs no
dedicated construct in the grammar. Comparisons such as $=$ and $\leq$
are recorded as $\emph{op}(\dataflowchain_1, \dataflowchain_2)$, and
boolean structure in filtering conditions is recorded by conjunction
and disjunction of dataflows. Grouping and ordering constructs,
\sqlgroupbycolor, \sqlorderbycolor, and \sqlpartitionbycolor, are
recorded directly, since they shape the computation in ways no other
construct captures.

\begin{figure}[!t]
\vspace{10pt}
\footnotesize 
\centering
\[
\arraycolsep=1.5pt\def\arraystretch{1.1}
\begin{array}{rrll}
&
\dataflowchain & 
\grammareq & 
\emph{an input table} \ | \ 
\emph{a column from an input table}  \\ 
& & \ \ \ | &  
\agg(\dataflowchain, \mydots, \dataflowchain) \ | \ 
\windowfunc \ | \ 
[\dataflowchain, \mydots, \dataflowchain] \ | \ 
\emph{op}(\dataflowchain, \dataflowchain) \ | \ 
\dataflowchain \land \dataflowchain \ | \ 
\dataflowchain \lor \dataflowchain \\ 
& & \ \ \ | &  
\sqlpartitionby(\dataflowchain) \ | \ 
\sqlorderby(\dataflowchain) \ | \ 
\sqlgroupby(\dataflowchain) \\ 
\end{array}
\]
\vspace{-12pt}
\caption{Dataflow language.}
\label{slab:fig:dataflowlanguage}
\vspace{-12pt}
\end{figure}

\begin{example}
Recall $\query_1$ from Table~\ref{slab:tab:Q1-Q2}, the human-written
query that, for each gender, computes a running total of scores by
day. Among the dataflows $\query_1$ manifests are:
{\small
\[
 \left \{
\text{\texttt{scores.score}, \sqlsumcolor(\texttt{scores.score}), $\texttt{scores.day} \leq \texttt{scores.day}$, $\sqlgroupbycolor \ \texttt{scores.gender}$, $\dotsb$ }
\right \}
\]
}
Reading them off: $\query_1$ draws data from the \texttt{score} column
of the \texttt{scores} table; it sums that column; an intermediate
step compares \texttt{day} values against each other; and it groups by
gender. Now consider $\query_2$ from the same table, the rewrite that
replaces $\query_1$'s self-join with a window function. Syntactically
the two queries share almost nothing, yet the essential computations
recur: $\query_2$ also draws on \texttt{scores.score}, also sums it,
and it organizes the computation around gender and day through its
\sqlpartitionbycolor and \sqlorderbycolor, constructs that differ from
$\query_1$'s \sqlgroupbycolor and day comparison in form but not in
role. Making this last correspondence precise requires one more
ingredient, alternative dataflows, introduced in
\cref{slab:sec:dataflow:extraction}; with it in place, the relation
between $\query_1$ and $\query_2$ lives entirely at the level of
dataflows, which is exactly the level at which a search can exploit
it.
\end{example}

\subsection{Extracting Dataflows}
\label{slab:sec:dataflow:extraction}

For dataflows to guide a search they must be computable, and cheaply,
since the search will consult them constantly. This subsection gives a
syntax-directed extraction procedure that computes the full set of
dataflows a query manifests in a single traversal of its structure.

The procedure is presented in \cref{slab:fig:dataflowextraction} as a
set of inference rules over judgments of the form
\[\centering
\calcalldataflowchains{\query_{\emph{ctx}}}{(\program)}{ \dataflowset }
\]
read: given the context $\query_{\emph{ctx}}$ in which the program
fragment $\program$ occurs, $\dataflowset$ is the set of all dataflows
manifested in $\program$. The fragment $\program$ ranges over whole
queries, filtering conditions, target lists, and expressions. The
context is what lets extraction see through a query's local
vocabulary: when $\program$ is, say, a target list,
$\query_{\emph{ctx}}$ is the query producing the table that
$\program$'s columns refer to, and the rules use it to chase every
column reference back to the input tables it ultimately draws from.

Two mutually recursive functions share the work. $\textsc{AllDfs}$
computes, for a fragment, the set of \emph{all} dataflows appearing
anywhere within it, components included. $\textsc{Df}$ computes the
\emph{single} dataflow denoted by the fragment itself, ignoring its
components. The distinction is what lets the rules record both facts
about a compound condition such as a disjunction: that each disjunct's
computations occur (via $\textsc{AllDfs}$), and that the disjunction
itself combines them (via $\textsc{Df}$).

Walking through the rules of \cref{slab:fig:dataflowextraction}: rule
(1) is the base case, extracting from a query that is just an input
table $\dbtable$ the singleton set containing $\dbtable$. Rules (2)
and (3) handle selection and join. Both take the union of the dataflow
sets of their constituents: for a selection, the source query, the
filtering condition, and the target list; for a join, the two sides
and the join condition. Neither construct mints a dataflow of its own;
selection is captured through composition, and the join logic is
carried by the join condition's dataflows. Rules (4) and (5) extract
from conditions. For a boolean combination, rule (4) collects all
dataflows of each operand condition and adds the dataflow of the
combination itself, obtained by applying the boolean operator to the
operands' $\textsc{Df}$ results; rule (5) does the corresponding work
for an atomic comparison over expressions. Rules (6) and (7) define
$\textsc{Df}$ on conditions, composing the operand dataflows under the
logical or comparison operator. Rules (8) through (13) ground
everything in expressions: a constant denotes itself (8); a column
$\col$ of an input table denotes the base dataflow $\dbtable.\col$
(9); a column reached through a join is resolved by recursing into
whichever side it comes from (10); a column reached through a
selection is resolved by locating its defining target in the target
list and recursing on that target (11); and aggregate applications
wrap the dataflows of their argument columns (12, 13).

\begin{figure*}[!t]
\footnotesize 
\resetmycounter{mycounter}
% NOTE (dissertation): content was sized for acmart's full page width; scaled to fit.
\begin{center}\resizebox{\textwidth}{!}{$\displaystyle
\arraycolsep=.5pt\def\arraystretch{1.1}
\begin{array}{l}
% first block 
\begin{array}{lll}
\fbox{\textsc{AllDfs} \text{algorithm that extracts all dataflows for queries}}
\\[5pt]
\begin{array}{cc}
(\showmycounter{mycounter}) & 
\irule{
\dbtable \text{ is an input table}
}{
\calcalldataflowchains{}{(\dbtable)}{ \{ \dbtable \} }
}
\end{array}
\begin{array}{cc}
(\showmycounter{mycounter}) & 
\irule{
\calcalldataflowchains{}{(\query)}{\dataflowset_1}  \ \ \ \ 
\calcalldataflowchains{\query}{(\sqlpredicate)}{\dataflowset_2} \ \ \ \ 
\calcalldataflowchains{\query}{(\targetlist)}{\dataflowset_3}
}{
\calcalldataflowchains{}{(\sqlselect \sqlspace \targetlist \sqlspace \sqlfrom \sqlspace \query \sqlspace \sqlwhere \sqlspace \sqlpredicate)}{\dataflowset_1 \cup \dataflowset_2 \cup \dataflowset_3}
}
\end{array}
\begin{array}{cc}
(\showmycounter{mycounter}) & 
\irule{
\calcalldataflowchains{}{(\query_1)}{\dataflowset_1} \ \ \ \ 
\calcalldataflowchains{}{(\query_2)}{\dataflowset_2} \ \ \ \ 
\calcalldataflowchains{\query_1 \sqlspace \sqljoin \sqlspace \query_2}{(\sqlpredicate)}{\dataflowset_3}
}{
\calcalldataflowchains{}{(\query_1 \sqlspace \sqljoin \sqlspace \query_2 \sqlspace \sqlon \sqlspace \sqlpredicate)}{ \dataflowset_1 \cup \dataflowset_2 \cup \dataflowset_3 } 
}
\end{array}
\end{array}
\\ \\ 
% second block 
\begin{array}{l}
\begin{array}{lll}
\fbox{\textsc{AllDfs} \text{algorithm that extracts all dataflows for conditions}}
\\[5pt]
\begin{array}{cc}
(\showmycounter{mycounter}) & 
\irule{
\calcalldataflowchains{\query}{(\sqlpredicate_i)}{\dataflowset_i} \ \ \ \ 
\calcdataflowchain{\query}{(\sqlpredicate_i)}{\dataflowchain_i} 
}{
\calcalldataflowchains{\query}{(\sqlpredicate_1 \ \emph{lop} \ \sqlpredicate_2)}{\dataflowset_1 \cup \dataflowset_2 \cup \{\emph{lop}(\dataflowchain_1, \dataflowchain_2) \}}
}
\end{array}
\begin{array}{cc}
(\showmycounter{mycounter}) & 
\irule{
\calcalldataflowchains{\query}{(\sqlexpr_i)}{\dataflowset_i} \ \ \ \ 
\calcdataflowchain{\query}{(\sqlexpr_i)}{\dataflowchain_i} 
}{
\calcalldataflowchains{\query}{(\sqlexpr_1 \emph{ op } \sqlexpr_2)}{\dataflowset_1 \cup \dataflowset_2 \cup \{ \emph{op}(\dataflowchain_1 , \dataflowchain_2) \}}
}
\end{array}
\\ \\ 
\fbox{\textsc{Df} \text{algorithm that extracts dataflows for conditions, expressions and targets}}
\\[5pt]
\begin{array}{cc}
(\showmycounter{mycounter}) & 
\irule{
\calcdataflowchain{\query}{(\sqlpredicate_i)}{\dataflowchain_i}
}{
\calcdataflowchain{\query}{(\sqlpredicate_1 \ \emph{lop} \  \sqlpredicate_2)}{ \emph{lop}(\dataflowchain_1, \dataflowchain_2) }
}
\end{array}
\begin{array}{cc}
(\showmycounter{mycounter}) & 
\irule{
\calcdataflowchain{\query}{(\sqlexpr_i)}{\dataflowchain_i}
}{
\calcdataflowchain{\query}{(\sqlexpr_1 \emph{ op } \sqlexpr_2)}{ \emph{op}(\dataflowchain_1 , \dataflowchain_2)  }
}
\end{array}
\begin{array}{cc}
\vspace{-5pt}
(\showmycounter{mycounter}) & 
\irule{
}{
\calcdataflowchain{\query}{(\emph{const})}{\emph{const}}
}
\end{array}
\begin{array}{cc}
(\showmycounter{mycounter}) & 
\irule{
\dbtable \text{ is an input table}
}{
\calcdataflowchain{\dbtable}{(\col)}{{T.}\col}
}
\end{array}
\begin{array}{cc}
(\showmycounter{mycounter}) & 
\irule{
\col \in \textsc{Columns}(\query_i) \ \ \ \ 
%\col \not\in \textsc{Columns}(\query_{2-i}) \ \ \ \ 
\calcdataflowchain{\query_i}{(\col)}{\dataflowchain_i}
}{
\calcdataflowchain{\query_1 \sqlspace \sqljoin \sqlspace \query_2 \sqlspace \sqlon \sqlspace \sqlpredicate}{(\col)}{ \dataflowchain_i } 
}
\end{array}
\\ \\ 
\begin{array}{cc}
(\showmycounter{mycounter}) & 
\irule{
(\target \sqlspace \sqlas \sqlspace \col) \in \targetlist \ \ \ \ 
\calcdataflowchain{\query}{(\target)}{\dataflowchain}
}{
\calcdataflowchain{\sqlselect \sqlspace \targetlist \sqlspace \sqlfrom \sqlspace \query \sqlspace \sqlwhere \sqlspace \sqlpredicate}{(\col)}{ \dataflowchain } 
}
\end{array}
\begin{array}{cc}
(\showmycounter{mycounter}) & 
\irule{
\calcdataflowchain{\query}{(\col)}{\dataflowchain}
}{
\calcdataflowchain{\query}{(\agg(\col))}{\agg(\dataflowchain)}
}
\end{array}
\begin{array}{cc}
(\showmycounter{mycounter}) & 
\irule{
\calcdataflowchain{\query}{(\col)}{\dataflowchain} 
}{
\calcdataflowchain{\query}{\big( \agg(\col) \sqlspace \sqlover( \sqlpartitionby \sqlspace \cols_1 \sqlspace \sqlorderby \sqlspace \cols_2 ) \big)}{ \agg(\dataflowchain) } 
}
\end{array}

\end{array}
\end{array}
\end{array}
$}\end{center}
\vspace{-10pt}
\caption{Inference rules that explain how dataflow extraction works for some key constructs in our query language.}
\label{slab:fig:dataflowextraction}
\vspace{-5pt}
\end{figure*}

\begin{example}
Consider the comparison
$$
\text{\texttt{s2.day <= s1.day}}
$$
inside $\query_1$ from Table~\ref{slab:tab:Q1-Q2}, where \texttt{s1}
and \texttt{s2} both alias the \texttt{scores} table. Extraction
produces the union of three sets: the dataflows of the left argument,
$\{\texttt{scores.day}\}$; the dataflows of the right argument,
likewise $\{\texttt{scores.day}\}$; and, from the comparison itself,
$\{\texttt{scores.day} \leq \texttt{scores.day}\}$. Note how the
context resolved both aliases to the same underlying input column: the
extracted dataflows say, in a syntax-independent way, that $\query_1$
reads \texttt{day} values and compares \texttt{day} values against
each other.
\end{example}

\head{Alternative dataflows}
The dataflow language, as presented so far, records operations in the
form the query expresses them, and this is occasionally too literal.
SQL offers distinct constructs that are not semantically equivalent in
isolation but can play corresponding roles in equivalent query
formulations. For example, a comparison
$\dataflowchain_1 \leq \dataflowchain_2$ and an ordering
$\sqlorderbycolor$ over the compared dataflows both record sensitivity
to the order of those values. Likewise, $\sqlgroupbycolor \
\dataflowchain$ and $\sqlpartitionbycolor \ \dataflowchain$ both
organize tuples by $\dataflowchain$, although one collapses each group
and the other retains its rows. Two equivalent queries can therefore
manifest related computations through literally different dataflow
terms; $\query_1$ and $\query_2$ from the running example are exactly
such a pair. We define selected pairs of role-corresponding dataflow
forms to be \emph{alternatives} of one another. This is a search
abstraction, not a claim that the underlying SQL constructs are
interchangeable. The ordering and grouping pairs above are two such
families, and our implementation recognizes further families of the
same character. When extracting the
dataflow set of the input query, the set is closed under alternatives:
whenever a dataflow is extracted, its alternative forms are added
alongside it. The closure is applied to the input query's set only.
Its purpose is to prevent a candidate from being treated as foreign
merely for expressing a corresponding role through a different
construct; candidates' own sets are left literal, so the closure
widens what counts as shared without manufacturing shared dataflows
out of nothing. With the closure in place, $\query_2$'s
$\sqlpartitionbycolor \ \texttt{scores.gender}$ registers as shared
with $\query_1$'s $\sqlgroupbycolor \ \texttt{scores.gender}$, and its
$\sqlorderbycolor \ \texttt{scores.day}$ as shared with $\query_1$'s
day comparison, so the two queries' kinship is fully visible to the
abstraction. Like the monotonicity observation below, the notion of
alternative dataflows was developed jointly with Rui Dong.

\head{A compositionality property}
The extraction rules only ever accumulate: the dataflow set of a
fragment always contains the dataflow sets of its components.
Consequently, whenever programs $\program_1, \mydots, \program_n$ are
composed into a larger program $\program$, the dataflow set of
$\program$ includes each component's set. This monotonicity is a
property of the abstraction itself, and the next section shows it is
what makes dataflows compatible with search: it guarantees that the
guidance signal degrades predictably as candidate programs grow, which
the synthesis algorithm exploits to prune entire regions of the space
at once. The observation that extraction is monotonic, and that this
enables efficient search, is mainly Rui Dong's contribution.

Dataflows, then, are precisely what the introduction asked for: a
semantic characterization of a query, mechanically computable from the
query alone, under which equivalent formulations are visibly related
and foreign computations are visibly foreign. What remains is to turn
this characterization into an engine, a way of steering an enumerative
synthesizer by how a candidate's dataflows relate to the input's. That
machinery, the contribution of Rui Dong, is the subject of the next
section.
